% ==========================================================
%  preamble/preamble.tex — การตั้งค่าเอกสารและแพ็กเกจหลัก
% ==========================================================

\usepackage{fontspec}
\usepackage{geometry}
\usepackage{graphicx}
\usepackage{booktabs}
\usepackage{longtable}
\usepackage{array}
\usepackage{multirow}
\usepackage{amsmath, amssymb}
\usepackage{caption}
\usepackage{float}
\usepackage{xcolor}
\usepackage{enumitem}
\usepackage{fancyhdr}
\usepackage{titlesec}
\usepackage{tabularx}
\usepackage{ragged2e}
\usepackage[hidelinks, bookmarks=true, bookmarksnumbered=true]{hyperref}

% ---------- ระยะขอบตามมาตรฐานวิทยานิพนธ์ ----------
\geometry{a4paper, top=3.0cm, bottom=2.5cm, left=3.5cm, right=2.5cm}

% ---------- เส้นทางค้นหารูปภาพ (Graphics Path) ----------
\graphicspath{{../figures/}{./figures/}}

% ---------- ฟอนต์ไทย: เลือกอัตโนมัติตามที่มีในเครื่อง ----------
% ลำดับความสำคัญ: TH Sarabun New (มาตรฐานราชการ) > Leelawadee UI > Angsana New > Tahoma
\newcommand{\THAIFONT}{}
\IfFontExistsTF{TH Sarabun New}{
  \renewcommand{\THAIFONT}{TH Sarabun New}
}{
  \IfFontExistsTF{Leelawadee UI}{
    \renewcommand{\THAIFONT}{Leelawadee UI}
  }{
    \IfFontExistsTF{Angsana New}{
      \renewcommand{\THAIFONT}{Angsana New}
    }{
      \renewcommand{\THAIFONT}{Tahoma}
    }
  }
}

% AutoFakeSlant/AutoFakeBold ช่วยให้ฟอนต์ไทยที่ไม่มีตัวเอียงยังแสดงผลได้
\setmainfont{\THAIFONT}[Ligatures=TeX, AutoFakeSlant=0.2, AutoFakeBold=1.5]
\setsansfont{\THAIFONT}[Ligatures=TeX, AutoFakeSlant=0.2, AutoFakeBold=1.5]
\newfontfamily\thaifont{\THAIFONT}[Ligatures=TeX]

% ฟอนต์ monospace สำหรับโค้ด (ใช้ฟอนต์ละตินเพื่อความคมชัด)
\IfFontExistsTF{Consolas}{
  \setmonofont{Consolas}[Scale=0.88]
}{
  \setmonofont{Courier New}[Scale=0.88]
}

% ---------- การตัดคำภาษาไทย ----------
% XeTeX ต้องใช้ locale "th" เพื่อตัดคำไทยได้ถูกต้อง
\XeTeXlinebreaklocale "th"
\XeTeXlinebreakskip = 0pt plus 0.1pt minus 0.05pt
\XeTeXlinebreakpenalty = 0
\sloppy
\emergencystretch=1em
\hbadness = 10000
\vbadness = 10000

% ---------- ระยะห่างบรรทัด ----------
\linespread{1.25}
\setlength{\parindent}{1.0cm}
\setlength{\parskip}{0.35em}

% ---------- รูปแบบหัวข้อ ----------
\titleformat{\chapter}[display]
  {\normalfont\huge\bfseries\centering}
  {\chaptertitlename\ \thechapter}{16pt}{\Huge}
\titlespacing*{\chapter}{0pt}{-20pt}{28pt}

\titleformat{\section}{\normalfont\Large\bfseries}{\thesection}{1em}{}
\titleformat{\subsection}{\normalfont\large\bfseries}{\thesubsection}{1em}{}
\titleformat{\subsubsection}{\normalfont\normalsize\bfseries}{\thesubsubsection}{1em}{}

% ---------- ชื่อภาษาไทยของส่วนต่าง ๆ ----------
\renewcommand{\contentsname}{สารบัญ}
\renewcommand{\listfigurename}{สารบัญภาพ}
\renewcommand{\listtablename}{สารบัญตาราง}
\renewcommand{\chaptername}{บทที่}
\renewcommand{\figurename}{ภาพที่}
\renewcommand{\tablename}{ตารางที่}
\renewcommand{\bibname}{บรรณานุกรม}
\renewcommand{\appendixname}{ภาคผนวก}

% ---------- รูปแบบคำบรรยายภาพและตาราง ----------
\captionsetup{font=normalsize, labelfont=bf, justification=centering,
              skip=6pt, labelsep=space}
\captionsetup[table]{position=top}
\captionsetup[figure]{position=bottom}

% ---------- หัวกระดาษและท้ายกระดาษ ----------
\pagestyle{fancy}
\fancyhf{}
\fancyfoot[C]{\thepage}
\renewcommand{\headrulewidth}{0pt}
\fancypagestyle{plain}{\fancyhf{}\fancyfoot[C]{\thepage}\renewcommand{\headrulewidth}{0pt}}
